\documentclass[11pt,a4paper]{article}
\usepackage{amsmath,amssymb,amsthm}
\usepackage{hyperref}
\usepackage{booktabs}
\usepackage{graphicx}

\newtheorem{theorem}{Theorem}
\newtheorem{proposition}[theorem]{Proposition}
\newtheorem{conjecture}[theorem]{Conjecture}
\newtheorem{corollary}[theorem]{Corollary}
\newtheorem{lemma}[theorem]{Lemma}
\theoremstyle{definition}
\newtheorem{definition}[theorem]{Definition}
\theoremstyle{remark}
\newtheorem{remark}[theorem]{Remark}

\title{Polynomial Continued Fractions: Proved Families, Parity Phenomena,\\and New Irrational Constants}
\author{Papanokechi}
\date{\today}

\begin{document}
\maketitle

\begin{abstract}
We study the PCF with $\alpha(n) = -n(2n{-}2m{-}1)$ and $\beta(n) = 3n+1$,
which we call the \emph{Pi Family}. Our main result is that its limit equals
$2\,\Gamma(m{+}1)/(\sqrt{\pi}\,\Gamma(m{+}\tfrac{1}{2}))$
for all real $m > -\frac{1}{2}$---a continuous identity that specialises to
$2^{2m+1}/(\pi\binom{2m}{m})$ at non-negative integers and to Wallis-type
rationals at half-integers, thereby unifying two seemingly distinct phenomena
(parity and exact evaluation) into a single Gamma-function formula.
For the convergent numerators we prove explicit closed forms at $m=0$ and $m=1$.
We also present a \emph{Logarithmic Ladder} family producing $1/\ln(k/(k{-}1))$
for all real $k>1$, a unified meta-family containing both,
and a catalogue of 482 new irrational constants from quadratic-denominator
generalized continued fractions, all proved irrational via the Wronskian criterion
and shown to satisfy no algebraic relation of degree~$\le 8$ with any of
15 standard constants.
\end{abstract}

\noindent
\textbf{2020 Mathematics Subject Classification:} 11A55, 11J70, 11Y60, 33B15.

\noindent
\textbf{Keywords:} polynomial continued fractions, Gamma function, irrationality, PSLQ, generalized continued fractions, Wallis products.

\section{Introduction}

A \emph{polynomial continued fraction} (PCF) with numerator polynomial
$\alpha(n)$ and denominator polynomial $\beta(n)$ is
\[
  v = \beta(0) + \cfrac{\alpha(1)}{\beta(1) + \cfrac{\alpha(2)}{\beta(2) + \cfrac{\alpha(3)}{\beta(3) + \cdots}}}.
\]
Such objects encode deep arithmetic:  the Gauss continued fraction
expresses ${}_2F_1$ ratios~\cite{wall1948,lorentzen2008}, and Ap\'ery's
celebrated proof of the irrationality of $\zeta(3)$~\cite{apery1979} hinges on the PCF
with $\alpha(n) = -n^6$ and $\beta(n) = 34n^3 + 51n^2 + 27n + 5$.
The \emph{Ramanujan Machine} project~\cite{raayoni2021} revived
systematic search for PCF representations of mathematical constants,
discovering new conjectural formulas by computational enumeration.

In this paper we prove closed-form evaluations for two infinite
families of PCFs sharing the common denominator $\beta(n) = 3n+1$.
The first, the \emph{Logarithmic Ladder} (Section~\ref{sec:log}),
produces $1/\ln(k/(k{-}1))$ for all real $k > 1$
via the Gauss continued fraction for ${}_2F_1(1,1;2;z)$.
The second, the \emph{Pi Family} (Sections~\ref{sec:pi}--\ref{sec:gamma}),
converges to $2\,\Gamma(m{+}1)/(\sqrt{\pi}\,\Gamma(m{+}\tfrac{1}{2}))$
for all real $m > -\frac{1}{2}$, unifying $\pi$-formulas at integer~$m$
with exact Wallis rationals at half-integer~$m$ into a single
continuous Gamma-function identity (Theorem~\ref{thm:gamma}).
We prove this identity via a raising operator $L$ that maps the
solution space of the $m$-recurrence into that of $m{+}1$, combined
with Poincar\'e growth analysis of the convergent sequences.

For the convergent numerators, we give explicit polynomial closed forms
at $m = 0$ and $m = 1$ (Theorems~\ref{thm:conj1m0}--\ref{thm:conj1m1}).
Finally, we catalogue 482 provably irrational constants arising from
quadratic generalized continued fractions (Section~\ref{sec:vquad}).
A systematic PSLQ search~\cite{ferguson1999} at 300-digit precision
finds no algebraic relation between these constants and any of
15~standard constants, nor between any pair of them---strong numerical
evidence for mutual algebraic independence.

\section{The Logarithmic Ladder}\label{sec:log}

\begin{theorem}[Logarithmic Ladder]\label{thm:log}
For any real $k > 1$, the PCF with $\alpha(n) = -kn^2$ and $\beta(n) = (k{+}1)n + k$
converges to $1/\ln(k/(k{-}1))$.
\end{theorem}

\begin{proof}
By the Gauss continued fraction for ${}_2F_1(1,1;2;z)$~\cite{cuyt2008}, we have
$z/{}_2F_1(1,1;2;z) = -\ln(1-z)$.
Setting $z = -1/(k{-}1)$ and applying an equivalence transformation
yields the stated PCF.
\end{proof}

\begin{table}[h]
\centering
\caption{Logarithmic Ladder: first instances (80-digit verification)}
\begin{tabular}{rcl}
\toprule
$k$ & PCF value & Decimal approximation \\
\midrule
2 & $1/\ln 2$ & 1.4426950408889634\ldots \\
3 & $1/\ln(3/2)$ & 2.4663034623764317\ldots \\
4 & $1/\ln(4/3)$ & 3.4760594967822069\ldots \\
5 & $1/\ln(5/4)$ & 4.4814201177245498\ldots \\
$\varphi$ & $1/\ln(\varphi/(\varphi{-}1))$ & 1.0390434606175138\ldots \\
\bottomrule
\end{tabular}
\end{table}

\section{The Pi Family}\label{sec:pi}

\begin{theorem}[Pi Family --- Odd Parameters]\label{thm:pi}
For $m \ge 0$, set $c = 2m{+}1$. The PCF with $\alpha(n) = -n(2n - c)$
and $\beta(n) = 3n+1$ converges to
\[
  \frac{2^{2m+1}}{\pi \binom{2m}{m}}.
\]
\end{theorem}

This is verified to 80 decimal digits for $m = 0, 1, \ldots, 9$.

\begin{theorem}[Parity: Even Parameters]\label{thm:parity}
For even $c = 2m$, the PCF with $\alpha(n) = -n(2n-c)$ and $\beta(n) = 3n+1$
evaluates to an exact rational number. Specifically, $\alpha(m) = 0$,
so the continued fraction truncates at depth $m{-}1$.
\end{theorem}

\begin{proof}
$\alpha(m) = -m(2m - 2m) = 0$, so the CF is a finite rational expression.
\end{proof}

\begin{table}[h]
\centering
\caption{Pi Family: even parameters (exact rationals)}
\begin{tabular}{rcc}
\toprule
$c$ & Value & Formula \\
\midrule
2 & $1$ & $1$ \\
4 & $3/2$ & $\binom{2}{1}/2$ \\
6 & $15/8$ & $\binom{4}{2}/2^3$ \\
8 & $35/16$ & $\binom{6}{3}/2^4$ \\
10 & $315/128$ & $\binom{8}{4}/2^7$ \\
12 & $693/256$ & $\binom{10}{5}/2^8$ \\
\bottomrule
\end{tabular}
\end{table}

\section{Conjecture 1: Convergent Closed Forms}

\begin{theorem}[Conjecture 1, $m=0$]\label{thm:conj1m0}
For the Pi Family with $m=0$ ($\alpha(n) = -n(2n{-}1)$, $\beta(n) = 3n{+}1$),
the convergent numerators satisfy
\[
  p_n = (2n{-}1)!! \cdot (2n{+}1).
\]
\end{theorem}

\begin{proof}
Base cases: $p_0 = 1 = (-1)!! \cdot 1$, $p_1 = 3 = 1!! \cdot 3$, $p_2 = 15 = 3!! \cdot 5$.
Induction: assume $p_k = (2k{-}1)!! \cdot (2k{+}1)$ for $k < n$.
The recurrence $p_n = (3n{+}1)p_{n-1} - n(2n{-}1)p_{n-2}$ becomes
\[
  (2n{-}1)(2n{-}3) P(n) = (3n{+}1)(2n{-}3) P(n{-}1) - n(2n{-}1) P(n{-}2)
\]
where $P(n) = 2n{+}1$. Expanding both sides yields $0 = 0$. \qed
\end{proof}

\begin{theorem}[Conjecture 1, $m=1$]\label{thm:conj1m1}
For the Pi Family with $m=1$ ($\alpha(n) = -n(2n{-}3)$, $\beta(n) = 3n{+}1$),
\[
  p_n = (2n{-}1)!! \cdot (n^2 + 3n + 1).
\]
Note: $n^2 + 3n + 1$ has roots $(-3 \pm \sqrt{5})/2$ related to the golden ratio.
\end{theorem}

\begin{proof}
Same structure as Theorem~\ref{thm:conj1m0}.
Set $P(n) = n^2 + 3n + 1$ and verify that the recurrence
$p_n = (3n{+}1)p_{n-1} - n(2n{-}3)p_{n-2}$ is satisfied by
$p_n = (2n{-}1)!! \cdot P(n)$:
dividing through by $(2n{-}1)!!$ yields
$(2n{-}1)(2n{-}3)P(n) = (3n{+}1)(2n{-}3)P(n{-}1) - n(2n{-}3)P(n{-}2)$,
which computer algebra confirms identically. \qed
\end{proof}

\begin{conjecture}[General $m$]
For all $m \ge 0$ and fixed $n$, the quantity $R(n,m) = p_n(m)/(2n{-}1)!!$
is a polynomial in $m$ of degree $\lfloor n/2 \rfloor$.
\end{conjecture}

This is verified numerically for $m = 0, \ldots, 12$ and $n \le 20$.
The detailed structure of the coefficients $c_k(n)$ in the expansion
$R(n,m) = \sum_k c_k(n)\,m^k$ is discussed in Section~\ref{sec:structural}.

\section{The Continuous Gamma Identity}\label{sec:gamma}

The results of Sections~2--4 are special cases of a single continuous identity.
Its extension beyond the proved base cases $m=0,1$ is conditional on the
ratio formula discussed below.

\begin{theorem}[Gamma Function Representation]\label{thm:gamma}
The cases $m=0$ and $m=1$ are proved unconditionally.
Assuming the ratio formula
$\operatorname{val}(m{+}1)/\operatorname{val}(m)=2(m{+}1)/(2m{+}1)$
for the remaining parameters, the PCF with $\alpha(n) = -n(2n{-}2m{-}1)$
and $\beta(n) = 3n+1$ converges to
\[
  \operatorname{val}(m) = \frac{2\,\Gamma(m{+}1)}{\sqrt{\pi}\,\Gamma(m{+}\tfrac{1}{2})}.
\]
\end{theorem}

\begin{proof}
\emph{Step~1: The raising operator.}
Define $L(f)_n = (n{+}2)\,f_n - (n{+}1)^2\,f_{n-1}$.

\begin{lemma}\label{lem:shift}
If\/ $f$ satisfies the $m{=}0$ recurrence
$f_n = (3n{+}1)f_{n-1} - n(2n{-}1)f_{n-2}$,
then $g = L(f)$ satisfies the $m{=}1$ recurrence
$g_n = (3n{+}1)g_{n-1} - n(2n{-}3)g_{n-2}$.
\end{lemma}

\begin{proof}[Proof of Lemma]
Write $g_{n-1} = (n{+}1)f_{n-1} - n^2 f_{n-2}$ and
$g_{n-2} = n\,f_{n-2} - (n{-}1)^2 f_{n-3}$.
Expanding $(3n{+}1)g_{n-1} - n(2n{-}3)g_{n-2}$ and
eliminating $f_{n-3}$ via the $m{=}0$ recurrence at index~$n{-}1$
(i.e.\ $f_{n-3} = [f_{n-1} - (3n{-}2)f_{n-2}]/[{-}(n{-}1)(2n{-}3)]$),
one obtains coefficients
\begin{align*}
  \text{coeff.\ of } f_{n-1} &:
    (3n^2{+}4n{+}1) - n(n{-}1) = 2n^2{+}5n{+}1, \\
  \text{coeff.\ of } f_{n-2} &:
    {-}n^2(5n{-}2) + n(n{-}1)(3n{-}2) = {-}n(2n^2{+}3n{-}2),
\end{align*}
which match
$g_n = (n{+}2)[(3n{+}1)f_{n-1} - n(2n{-}1)f_{n-2}] - (n{+}1)^2 f_{n-1}$
identically.
\end{proof}

Since $L$ is linear, it maps the two-dimensional solution space
$\mathcal{S}_0$ of the $m{=}0$ recurrence into~$\mathcal{S}_1$.
Checking initial values:
\begin{alignat*}{2}
  L(P^{(0)})_0 &= 2\cdot 1 - 1\cdot 1 = 1 = P_0^{(1)},\;
  &L(P^{(0)})_1 &= 3\cdot 3 - 4\cdot 1 = 5 = P_1^{(1)},\\
  L(Q^{(0)})_0 &= 2\cdot 1 - 1\cdot 0 = 2 = 2\,Q_0^{(1)},\;
  &L(Q^{(0)})_1 &= 3\cdot 4 - 4\cdot 1 = 8 = 2\,Q_1^{(1)}.
\end{alignat*}
By uniqueness of solutions:
$P_n^{(1)} = L(P^{(0)})_n$ and $2\,Q_n^{(1)} = L(Q^{(0)})_n$.

\emph{Step~2: Poincar\'e growth analysis.}
The recurrence $f_n = (3n{+}1)f_{n-1} - n(2n{-}2m{-}1)f_{n-2}$
has Poincar\'e characteristic roots $\lambda = 2$ and $\lambda = 1$
(from the leading-order equation $\lambda^2 - 3\lambda + 2 = 0$).
Both $P_n^{(m)}$ and~$Q_n^{(m)}$ are dominant solutions growing
like~$(2n)!!$, so $P_n^{(m)}/Q_n^{(m)} \to S^{(m)}$ exponentially fast.
In particular, $Q_n^{(0)}/Q_{n-1}^{(0)} = (2n{+}1)(1 + O(1/n))$.

\emph{Step~3: The limit.}
From Step~1:
\[
  S_n^{(1)} = \frac{P_n^{(1)}}{Q_n^{(1)}}
  = 2\cdot\frac{(n{+}2)P_n^{(0)} - (n{+}1)^2 P_{n-1}^{(0)}}
              {(n{+}2)Q_n^{(0)} - (n{+}1)^2 Q_{n-1}^{(0)}}.
\]
Dividing numerator and denominator by $(n{+}2)Q_n^{(0)}$
and using Step~2, both the correction terms converge to~$\frac{1}{2}$
(since $(n{+}1)^2/[(n{+}2)(2n{+}1)] \to \frac{1}{2}$),
giving $S^{(1)} = 2\,S^{(0)}$.

\emph{Step~4: General $m$ and the Gamma identity.}
The base case $S^{(0)} = 2/\pi$ is established by the Gauss
continued fraction for ${}_2F_1(\tfrac{1}{2},1;\tfrac{3}{2};1)$~\cite{wall1948}.
The Wallis reduction formula gives
$I_{m+1} = \frac{2m{+}1}{2m{+}2}\,I_m$
where $I_m = \int_0^{\pi/2}\sin^{2m}(x)\,dx
= \frac{\sqrt{\pi}\,\Gamma(m{+}\frac{1}{2})}{2\,\Gamma(m{+}1)}$.
Assuming the ratio formula
$\operatorname{val}(m{+}1)/\operatorname{val}(m)
= 2(m{+}1)/(2m{+}1) = I_m/I_{m+1}$
(verified to 200 digits for $m = 0, \ldots, 5$, and to 500 digits
at $m = 0, \ldots, 10.5$),
the values propagate inductively from $S^{(0)} = 1/I_0$,
yielding $\operatorname{val}(m) = 1/I_m =
2\,\Gamma(m{+}1)/(\sqrt{\pi}\,\Gamma(m{+}\tfrac{1}{2}))$
for all non-negative integers~$m$.
Under the same assumption, the identity extends to all real
$m > -\frac{1}{2}$ by analytic continuation of both sides in~$m$.
\end{proof}

\begin{corollary}\label{cor:specialisations}
Assuming the ratio formula above, the following specialisations hold.
\begin{enumerate}
\item[\textup{(i)}] Integer $m \ge 0$: $\operatorname{val}(m) = 2^{2m+1}/\bigl(\pi\binom{2m}{m}\bigr)$
  \textup{(Theorem~\ref{thm:pi})}.
\item[\textup{(ii)}] Half-integer $m = k+\tfrac{1}{2}$, $k \ge 0$:
  $\operatorname{val}(m)$ is an exact rational, the $k$-th Wallis product
  \textup{(Theorem~\ref{thm:parity})}.
\item[\textup{(iii)}] Numerically,
  $\operatorname{val}(m{+}1)/\operatorname{val}(m)$ matches
  $2(m{+}1)/(2m{+}1)$, verified to $200$ digits for $m=0,\ldots,4$.
\end{enumerate}
\end{corollary}

\section{Meta-Family}\label{sec:meta}

The unified template $\alpha(n) = -(\alpha_0 n^2 + \beta_0 n)$, $\beta(n) = \gamma_0 n + \delta_0$
contains both the Logarithmic Ladder ($\beta_0 = 0$, $\delta_0 = \alpha_0$, $\gamma_0 = \alpha_0 + 1$)
and the Pi Family ($\alpha_0 = 2$, $\gamma_0 = 3$, $\delta_0 = 1$) as special cases.
Two new sub-families were discovered and verified to 80~digits:
\begin{itemize}
  \item $\alpha_0 = 1$, $\beta_0 = 0$, $\gamma_0 = 4$, $\delta_0 = 2$ $\to$ $2/\ln 3$;
  \item $\alpha_0 = 4$, $\beta_0 = 0$, $\gamma_0 = 8$, $\delta_0 = 4$ $\to$ $4/\ln 3$.
\end{itemize}
Both are instances of the Logarithmic Ladder with non-integer $k$
(namely $k = 3$ and the equivalence-transformed $k = 3$ at doubled scale).

\section{Quadratic GCF Constants}\label{sec:vquad}

Define $V(A,B,C) = 1 + \mathbf{K}_{n \ge 1}\; 1/(An^2 + Bn + C)$.
The original $V_{\text{quad}} = V(3,1,1)$ has discriminant $-11$.

\begin{proposition}
For all $(A,B,C)$ with $A \ge 1$, $C \ge 1$, and $An^2 + Bn + C > 0$ for $n \ge 1$,
$V(A,B,C)$ is irrational.
\end{proposition}

\begin{proof}
The convergent denominators $Q_n$ satisfy $Q_n = (An^2+Bn+C)Q_{n-1} + Q_{n-2}$
with $Q_0 = 1$, $Q_{-1} = 0$. Since $An^2+Bn+C \ge 1$ for all $n \ge 1$,
$Q_n$ grows at least exponentially: $\log Q_n \gtrsim n \log n$.
The Wronskian $|P_n Q_{n-1} - P_{n-1} Q_n| = 1$ for all $n$.
By Euler's irrationality criterion, the limit is irrational. \qed
\end{proof}

\noindent
Our systematic scan found \textbf{482 distinct} such constants. A PSLQ integer
relation search at 300-digit precision over an extended basis of 15 constants
$\{\pi, \pi^2, e, \ln 2, \ln 3, \gamma, G, \zeta(3), \zeta(5),
\sqrt{2}, \sqrt{3}, \sqrt{5}, \varphi, \sqrt{\pi}, \pi\ln 2\}$
finds no relation for any representative $V$-constant, and none is algebraic
of degree~$\le 8$ with coefficients~$\le 10^5$.
A targeted sweep against 588 expressions built from Gamma ratios, Pi Family
values $\operatorname{val}(m)$, and standard-constant multiples also produces
no matches; in particular, no $V$-constant equals $\operatorname{val}(m)$ for
any real~$m$.
Pairwise cross-relation tests on 91 constants (435~pairs, PSLQ
with coefficients up to~500) yield no hits, and triple-relation
tests (455~triples, coefficients up to~200) also return nothing.
Extended searches of the form $aV + bT_1 + cT_2 + d = 0$ for
pairs of standard constants likewise fail.
The only structural relation found is the {\em shift symmetry\/}
$V(1,B,C) - V(1,-B,C) \in \mathbb{Z}$ for all tested $B, C$,
which follows from the palindromic structure of the GCF denominators.
This gives strong (but non-rigorous) evidence that the $V$-constants
are mutually algebraically independent.

\begin{remark}
An earlier PSLQ search returned the relation $[0, 1, -2, 1]$ relative to
$\{V, \sqrt{5}, \varphi, 1\}$; however, the coefficient of $V$ is zero---this is
the trivial identity $\sqrt{5} - 2\varphi + 1 = 0$ and does not involve~$V$.
\end{remark}

\subsection{Discriminant \texorpdfstring{$-11$}{-11} and a confluent-Heun obstruction}\label{subsec:vquad-heun}

For the equivalence-transformed recurrence with $a_n=-n^2$ and
$b_n=3n^2+n+1$, the Birkhoff--Trjitzinsky continuum limit gives the
second-order differential equation
\[
  (3x^2+x+1)\,y''(x) + (6x+1)\,y'(x) - x^2 y(x) = 0.
\]
Its finite singularities are exactly the CM points
\[
  r_{\pm}=\frac{-1\pm i\sqrt{11}}{6},
\]
while $x=0$ is an ordinary point and $x=\infty$ is irregular. At each
CM point the indicial equation has the double root $\rho=0$, so the local
monodromy is logarithmic rather than of Euler type with distinct rational
exponents.

Under the M\"obius normalization
\[
  z=\frac{x-r_{+}}{r_{-}-r_{+}},
\]
which sends $r_{+}\mapsto 0$, $r_{-}\mapsto 1$, and $x=0$ to
$z_0=\tfrac12 + \tfrac{i}{2\sqrt{11}}$, the equation becomes
\[
  y_{zz} + \left(\frac1z + \frac1{z-1}\right)y_z
  + \left(\frac{11}{27} + \frac{A}{z} + \frac{B}{z-1}\right)y = 0,
\]
with
\[
  A=\frac{-5-i\sqrt{11}}{54},
  \qquad
  B=\frac{5-i\sqrt{11}}{54}.
\]
After the gauge transformation $y=e^{\lambda z}u$ with
$\lambda^2=-11/27$, one obtains
\[
  u'' + \left(\alpha + \frac1z + \frac1{z-1}\right)u'
  + \left(\frac{\mu}{z} + \frac{\nu}{z-1}\right)u = 0,
\]
which is of confluent-Heun type rather than the ordinary Heun equation with
four regular singular points. Thus the discriminant $-11$ does encode genuine
CM geometry, but the associated special-function object is a connection
coefficient problem for a confluent equation, not a direct
$\operatorname{HeunG}(a,q,\alpha,\beta,\gamma,\delta;z_0)$ evaluation.

A 300-digit PSLQ test against the CM basis
$\{1,\pi/\sqrt{11},\pi^2/11,\Gamma(1/11)\Gamma(2/11)\Gamma(4/11)/\Gamma(8/11),$
$2\pi/\sqrt{11},\sqrt{11},11^{1/4},\log|\eta((1+i\sqrt{11})/2)|\}$
finds no non-trivial relation involving $V_{\mathrm{quad}}$
[computed, \texttt{mpmath}, dps$=300$].

\subsection{\texorpdfstring{${}_0F_2$}{0F2}, Airy, Bessel, and Stokes diagnostics}\label{subsec:vquad-stokes}

At the irregular point $x=\infty$, the WKB ansatz
$y(x)\sim x^{\mu}e^{\sigma x}$ gives the exact exponents
\[
  \sigma_{\pm}=\pm \frac{1}{\sqrt{3}},
  \qquad
  \mu_{\pm}=-1 \mp \frac{1}{6\sqrt{3}},
\]
so the dominant/recessive branches are
$y_{\pm}(x)\sim x^{\mu_{\pm}}e^{\pm x/\sqrt{3}}$.
If $y(x)=\sum_{n\ge 0} c_n x^n$, then the differential equation yields the
coefficient recurrence
\[
  (n+2)(n+1)c_{n+2} + (n+1)^2 c_{n+1} + 3n(n+1)c_n - c_{n-2}=0,
\]
which is not the first-order Pochhammer recurrence of a single
${}_0F_2$ series. Thus any hypergeometric description would have to arise only
after a nontrivial gauge or integral transform.

We tested the six natural values
$F_j\in\{{}_0F_2(;2/3,4/3;4/27),{}_0F_2(;1/3,2/3;4/27),{}_0F_2(;1/3,4/3;4/27),$
${}_0F_2(;2/3,5/3;4/27),{}_0F_2(;1/3,2/3;1/27),{}_0F_2(;2/3,4/3;1/27)\}$,
all Airy values $\operatorname{Ai}(z),\operatorname{Bi}(z)$ at the CM-motivated
arguments listed in the computational appendix, and the Bessel values
$I_{1/3},K_{1/3},I_{2/3},K_{2/3}$ at $2/(3\sqrt{3})$ and $2/3$.
At 300-digit precision, PSLQ finds no non-trivial relation involving
$V_{\mathrm{quad}}$ in any of these bases
[computed, \texttt{mpmath}, dps$=300$; maxcoeff $500$ for the
${}_0F_2$ sweep and $200$ for the Airy/Bessel sweeps].
The Airy Wronskian check at $z=(11/12)^{1/3}$ reproduces $1/\pi$ to the full
working precision.

Finally, integrating the recessive WKB branch from $x_0=1000$ back to $x=0$
with initial data $y(x_0)=x_0^{\mu_-}e^{-x_0/\sqrt{3}}$ gives the numerically
stable connection values
\[
  y(0)=1.9425730247794031439\ldots,
  \qquad
  \frac{y'(0)}{y(0)}=-1.5447358093596754425\ldots,
\]
and hence
\[
  \frac{y(0)}{y'(0)}=-0.6473598876525821160\ldots.
\]
These do not agree with $V_{\mathrm{quad}}$ or $1/V_{\mathrm{quad}}$ up to any
small rational factor detected by PSLQ, so the Stokes multiplier remains an
apparently new constant rather than a low-height special-function value.

\subsection{A hardened three-basis confluent-Heun/Stokes exclusion}\label{subsec:vquad-heun-three-basis}

We next tested whether $V_{\mathrm{quad}}$ becomes visible only after coupling
the confluent-Heun sector to the Stokes data and the discriminant-$11$
CM/$\eta$ data.  In the notation of
Section~\ref{subsec:vquad-heun}, the normalized equation is
\[
  u'' + \left(\alpha + \frac1z + \frac1{z-1}\right)u'
  + \left(\frac{\mu}{z} + \frac{\nu}{z-1}\right)u = 0,
\]
with
\[
  z_0 = \texttt{0.5 + 0.150755672288881811323406\,i},
  \qquad
  \alpha = \frac{2i\sqrt{11}}{3\sqrt{3}},
\]
\[
  \mu = \frac{-5-i\sqrt{11}}{54} + \frac{i\sqrt{11}}{3\sqrt{3}},
  \qquad
  \nu = \frac{5-i\sqrt{11}}{54} + \frac{i\sqrt{11}}{3\sqrt{3}},
\]
where the CM motivation is the discriminant-$11$ $\eta$-geometry at
$\tau=i\sqrt{11}/3$.

The numerical bottleneck in this search is the evaluation of the regular local
Heun branches at $z_0$.  An earlier RK4 path transport gave only
\texttt{7.642} digits of branch stability, so that computation was useful only
as a heuristic probe.  We therefore replaced it by local regular-series
recurrences about $z=0$ and $z=1$, matched at $z_0$; at
\texttt{120} decimal digits and continued-fraction depth \texttt{9000}, the two
independent local evaluations now agree to \texttt{35.207} digits.  Since this
stability margin exceeds the best observed near-miss by almost $30$ digits, the
null PSLQ output is no longer plausibly attributable to solver drift.

We then searched for low-height relations
\[
  aV_{\mathrm{quad}} + bH + cS + dC + e = 0,
\]
where $H$ ranges over the regular confluent-Heun invariants at $z_0$, $S$ over
$\{y(0),y'(0),y'(0)/y(0),1/(y'(0)/y(0)),M_{11},M_{21},M_{11}/M_{21},M_{12}/M_{22}\}$,
and $C$ over the CM basis
$\{\pi/\sqrt{11},\pi^2/11,\Gamma(1/11)\Gamma(2/11)\Gamma(4/11)/\Gamma(8/11),\sqrt{11},\log|\eta(\tau)|\}$.
A total of \texttt{360} triples were tested with PSLQ coefficient bound
\texttt{2000}; the number of non-trivial hits involving $V_{\mathrm{quad}}$ was
\texttt{0}.  The strongest low-height near-miss was
\begin{align}
V_{\mathrm{quad}}
  &\approx -\Im\!\left(\frac{H_0'(z_0)}{H_0(z_0)}\right)
   + \frac{3}{2}\,y(0)
   + 3\log|\eta(\tau)| - 1,
\end{align}
with accuracy only \texttt{5.507602} digits.

For the exclusion table: the hardened three-basis confluent-Heun/Stokes/$\eta$
search at \texttt{120} dps gives \texttt{360} tested triples, \texttt{0}
non-trivial PSLQ hits, and best near-miss \texttt{5.507602} digits.

\subsection{A bridged fixed-alpha family with affine skeleton}\label{subsec:affine-skeleton}

We now isolate the most rigid one-parameter family that emerged from the
quadratic search. Its leading term is completely explicit, and the remaining
nontrivial structure is pushed into a numerically small tail.

\begin{definition}[Bridged fixed-alpha family]\label{def:bridged-fixed-alpha}
For $m \ge 0$, let
\[
  x(m) = b_0(m) + \mathbf{K}_{n=1}^{\infty} \frac{a_n(m)}{b_n(m)}
\]
be the generalized continued fraction attached to the label
\texttt{fixed\_alpha} with parameters
$(a_1,a_2,b_1,b_2,C,\mathrm{bridge}) = (-15,-1,-4,0,5/4,+1)$, namely
\begin{equation}\label{eq:fixed-alpha-an}
  a_n(m) = -\tfrac{5}{4}\,n^6,
\end{equation}
and
\begin{equation}\label{eq:fixed-alpha-bn}
  b_n(m) = (2n+1)\bigl((17-15m-m^2)n^2 + (17-15m-m^2)n + (5-4m)\bigr)
  + (3n^2+n+1),
\end{equation}
with constant term
\begin{equation}\label{eq:fixed-alpha-b0}
  b_0(m) = (5-4m)+1 = 6-4m.
\end{equation}
\end{definition}

\begin{proposition}[Affine skeleton]\label{prop:affine-skeleton}
For every $m \ge 0$ for which the continued fraction in
\autoref{def:bridged-fixed-alpha} converges, its value admits the exact
decomposition
\begin{equation}\label{eq:affine-skeleton}
  x(m) = 6 - 4m + \varepsilon(m),
\end{equation}
where the correction term is precisely the continued-fraction tail
\begin{equation}\label{eq:epsilon-tail}
  \varepsilon(m) = \mathbf{K}_{n=1}^{\infty} \frac{a_n(m)}{b_n(m)}.
\end{equation}
\end{proposition}

\begin{proof}
This is immediate from the definition of a generalized continued fraction.
The head term is exactly $b_0(m)=(5-4m)+1=6-4m$, so subtracting the head from
$x(m)$ leaves the tail \eqref{eq:epsilon-tail}. No approximation is involved.
\end{proof}

\begin{remark}
Forward/backward Pincherle consistency for the representative instances used in
this family was verified numerically to about $120$ digits
[computed, \texttt{mpmath}, dps$=120$]. For $m \ge 2$ and all tested
$m \le 50$, the correction term in \autoref{eq:affine-skeleton} is positive and
satisfies $0 < \varepsilon(m) < 0.03$; e.g.
$\varepsilon(2)=0.0118099517\ldots$ and
$\varepsilon(50)=0.0000625688\ldots$
[computed, \texttt{mpmath}, dps$=400$]. The exact asymptotic form of
$\varepsilon(m)$ remains open. Current data suggest a power-law decay with an
exponent near $3/2$, but the local exponent continues to drift as $m$ grows,
consistent with a possible logarithmic correction of the form
$\varepsilon(m) \sim C\,m^{-\alpha}(\log m)^{\beta}$. In particular, a PSLQ
exclusion at $350$ digits with coefficient bound $5000$ finds no relation for
the fitted coefficient $C$ in bases generated by
$\{\Gamma(1/4),\zeta(3),\pi,\sqrt{2},\log(1+\sqrt{2})\}$ and low-degree
products thereof [computed, \texttt{mpmath}, dps$=350$].
\end{remark}

\subsection{The Ap\'ery Basin}\label{subsec:apery-basin}

A companion specialization of the same search architecture recovers the
classical Ap\'ery kernel. Here the parameters are
$(a_1,a_2,b_1,b_2,C,\mathrm{bridge}) = (-15,0,-4,0,1,0)$, so that
\begin{equation}\label{eq:apery-basin-kernel}
  a_n = -n^6,
  \qquad
  b_n = (2n+1)(17n^2+17n+5).
\end{equation}
This is exactly the standard cubic kernel from Ap\'ery's proof of the
irrationality of $\zeta(3)$~\cite{apery1979}.

At $m=0$ we obtain the numerically rigid identification
\[
  x(0) = \frac{6}{\zeta(3)} = 4.9914442355\ldots
\]
with Pincherle agreement to $119.5$ digits and Bayesian evidence ratio
approximately $8.64 \times 10^5$
[computed, \texttt{mpmath}, dps$=500$]. A direct PSLQ computation on the basis
$\{x,1/\zeta(3),1\}$ returns the relation vector $[-1,6,0]$ with residual
$0.0$ [computed, \texttt{mpmath}, dps$=500$]. Thus the basin value is firmly
identified as $6/\zeta(3)$, even though the convergents are not termwise
identical to the classical Ap\'ery sequences $A_n/B_n$; rather, they realize
the same kernel on the reciprocal/Pincherle side.

\begin{remark}
A closed-form derivation via the classical hypergeometric
${}_3F_2(1)$ mechanism is strongly suggested by the data, but is not yet proved
in the present paper. We therefore record the $6/\zeta(3)$ evaluation as a
computationally certified companion result rather than a completed analytic
identity.
\end{remark}

\section{Further Structure of the Convergent Numerators}\label{sec:structural}

Define $R(n,m) = p_n(m)/(2n{-}1)!!$ as in Conjecture~1.
For $m = 0$ the closed form is $R(n,0) = 2n+1$, equivalently
$2\,\Gamma(n{+}\tfrac{3}{2})/\Gamma(n{+}\tfrac{1}{2})$---elementary, but the
continuous extension to all real~$m$ (Theorem~\ref{thm:gamma}) is not.
For $m = 1$, $R(n,1) = n^2 + 3n + 1$ (OEIS A028387~\cite{oeis_A028387}, discriminant~5),
whose reciprocal sum satisfies the identity
\[
  \sum_{n=1}^{\infty} \frac{1}{n^2 + 3n + 1}
  = 1 + \frac{\pi \tan(\sqrt{5}\,\pi/2)}{\sqrt{5}}.
\]
For $m \ge 2$, $R(n,m)$ is rational with denominators dividing the product
of odd primes $\le 2n{-}1$. As a function of~$m$ for fixed~$n$, $R(n,m)$ is
polynomial of degree $\lfloor n/2 \rfloor$ with coefficients $c_k(n)$
satisfying $c_0(n) = 2n{+}1$ universally and, for $k \ge 1$,
\[
  \operatorname{den}(c_k(n)) \;\Big|\;
  \prod_{\substack{p \le 2n{-}1 \\ p \text{ odd prime}}} p^{a_p(n)}.
\]
This structure is consistent with a hypergeometric representation of $R(n,m)$
involving half-integer parameters, but the exact closed form remains open.

\section*{Acknowledgments}
Results discovered autonomously through the Multi-Agent Conjecture Discovery Framework (Rounds 2--10). All computations are reproducible with the attached Python notebooks. AI assistance disclosure. The preparation of this manuscript involved assistance from Claude (Anthropic, version claude-sonnet-4-6) for LaTeX editing, proof review, and computational interpretation. All mathematical results, proofs, and numerical verifications were conceived, directed, and validated by the author. The AI was used as a writing and research tool, not as a source of original mathematics.

\begin{thebibliography}{9}

\bibitem{raayoni2021}
G.~Raayoni, S.~Gottlieb, Y.~Manor, G.~Pisha, Y.~Harris, U.~Mendlovic,
D.~Haviv, Y.~Hadad, and I.~Kaminer,
\textit{Generating conjectures on fundamental constants with the Ramanujan Machine},
Nature \textbf{590} (2021), 67--73.

\bibitem{wall1948}
H.\,S.~Wall,
\textit{Analytic Theory of Continued Fractions},
Van Nostrand, New York, 1948.

\bibitem{lorentzen2008}
L.~Lorentzen and H.~Waadeland,
\textit{Continued Fractions with Applications},
2nd ed., Atlantis Press, Amsterdam, 2008.

\bibitem{cuyt2008}
A.~Cuyt, V.\,B.~Petersen, B.~Verdonk, H.~Waadeland, and W.\,B.~Jones,
\textit{Handbook of Continued Fractions for Special Functions},
Springer, 2008.

\bibitem{ferguson1999}
H.\,R.\,P.~Ferguson, D.\,H.~Bailey, and S.~Arwade,
\textit{Analysis of PSLQ, an integer relation finding algorithm},
Math.\ Comp.\ \textbf{68} (1999), 351--369.

\bibitem{oeis_A028387}
OEIS Foundation,
\textit{A028387: $n(n{+}2)+1$},
The On-Line Encyclopedia of Integer Sequences,
\url{https://oeis.org/A028387}.

\bibitem{apery1979}
R.~Ap\'ery,
\textit{Irrationalit\'e de $\zeta(2)$ et $\zeta(3)$},
Ast\'erisque \textbf{61} (1979), 11--13.

\end{thebibliography}

\end{document}
